\documentclass[11pt,a4paper]{article}
\usepackage[utf8]{inputenc}
\usepackage[margin=1in]{geometry}
\usepackage{amsmath, amssymb}
\usepackage{booktabs}
\usepackage{tikz}
\usetikzlibrary{arrows.meta, angles, quotes, calc}
\usepackage{titlesec}
\titlespacing*{\section}{0pt}{1.4ex plus .2ex minus .2ex}{0.8ex plus .1ex}
\setlength{\abovedisplayskip}{7pt}
\setlength{\belowdisplayskip}{7pt}

\title{\texttt{IKCompiler.calc\_shoulder\_and\_elbow\_angles()}}
\author{}
\date{}

\begin{document}

\maketitle
\vspace{-5.5em}
\small

\begin{center}
\resizebox{0.56\textwidth}{!}{%
\begin{tikzpicture}[
    >=Stealth,
    pt/.style={circle, fill, inner sep=1.3pt},
    lbl/.style={font=\small},
    ra/.style={gray!20!black}
]
    \draw[gray!35, thin] (-0.8,-0.6) -- (7.0,-0.6);
    \draw[gray!35, thin] (-0.8,-0.6) -- (-0.8,4.2);

    \coordinate (S) at (0, 0);
    \coordinate (W) at (4.14, 1.51);
    \coordinate (E) at (0.81, 2.89);
    \coordinate (Ref) at (5.6, 0);

    % horizontal reference (world zero angle) and auxiliary chord
    \draw[dashed, gray!55] (S) -- (Ref);
    \draw[dashed, gray!70] (S) -- (W) node[midway, below right=-2pt, lbl, gray!50!black] {$c$};

    % physical links
    \draw[thick, blue!60!black] (S) -- (E) node[midway, left=3pt, lbl, blue!60!black] {$a$};
    \draw[thick, red!70!black] (E) -- (W) node[midway, above right=-2pt, lbl, red!70!black] {$b$};

    \node[pt, violet!70!black] at (S) {};
    \node[lbl, below=2pt] at (S) {$\mathbf s$};
    \node[pt, teal!70!black] at (E) {};
    \node[lbl, above left=1pt] at (E) {$E$};
    \node[pt, red!70!black] at (W) {};
    \node[lbl, right=2pt] at (W) {$\mathbf w$};

    % elevation angle epsilon: horizontal ref to chord c
    \pic[draw, violet!60!black, angle radius=7mm, angle eccentricity=1.5, "$\varepsilon$"] {angle = Ref--S--W};

    % gamma: chord c to link a, at S
    \pic[draw, orange!80!black, angle radius=11mm, angle eccentricity=1.4, "$\gamma$"] {angle = W--S--E};

    % delta: elbow bend angle at E
    \pic[draw, teal!70!black, angle radius=8mm, angle eccentricity=1.5, "$\delta$"] {angle = S--E--W};

\end{tikzpicture}}
\end{center}

\noindent
Physical links $\mathbf s E$ (length $a$) and $E\mathbf w$ (length $b$), solid. Auxiliary
straight-line chord $\mathbf s \mathbf w$ (length $c$) and world-horizontal reference, dashed.
$\gamma,\delta$ solved by the law of cosines on triangle $\mathbf s E\mathbf w$;
$\varepsilon$ is the chord's own elevation, so $\gamma+\varepsilon$ is link $a$'s elevation
in the world frame.

\section{Setup}

\begin{align*}
\mathbf s &= \texttt{shoulder\_servo\_position}, \quad
\mathbf w = \texttt{wrist\_position} &&\in \mathbb R^3\\
a &= \texttt{shoulder\_arm\_length}, \quad
b = \texttt{elbow\_arm\_length}
\end{align*}

$E$ is the (unknown, never explicitly computed) elbow position. It is fixed only implicitly,
as the point at distance $a$ from $\mathbf s$ and distance $b$ from $\mathbf w$, in the vertical
plane containing $\mathbf s$, $\mathbf w$, and the $z$-axis.

\section{Elevation Angle}

\begin{align}
\mathbf d &= \mathbf w - \mathbf s, \qquad
r = \sqrt{d_x^2 + d_y^2} \\
\varepsilon &= \operatorname{atan}\!\left(\frac{d_z}{r}\right)
\end{align}
the angle the chord $\mathbf s\mathbf w$ makes above the world-horizontal plane (\texttt{atan}
suffices in place of \texttt{atan2} since $r \ge 0$ always).

\begin{equation}
c = \lVert \mathbf d \rVert = \sqrt{r^2 + d_z^2}
\end{equation}

\section{Law of Cosines}

Triangle $\mathbf s E \mathbf w$ has sides $a = \mathbf s E$, $b = E\mathbf w$,
$c = \mathbf s \mathbf w$. The law of cosines gives its interior angles at $\mathbf s$ and $E$:
\begin{align}
\gamma &= \angle(\mathbf s) = \arccos\!\left(\frac{a^2 + c^2 - b^2}{2ac}\right)
&&\text{angle between } \mathbf s\mathbf w \text{ and } \mathbf s E \\
\delta &= \angle(E) = \arccos\!\left(\frac{a^2 + b^2 - c^2}{2ab}\right)
&&\text{angle between } E\mathbf s \text{ and } E\mathbf w
\end{align}
$\delta = 180^\circ$ is the fully-extended (straight) configuration $\mathbf s, E, \mathbf w$
collinear -- the reference the servo-space elbow command is built against
(\texttt{elbow\_logical\_zero + (180 - elbow\_angle)}).

\section{World-Frame Shoulder Angle}

$\gamma$ is measured from the chord $\mathbf s\mathbf w$, not from horizontal. Since
$\varepsilon$ is exactly the chord's own elevation, adding it rotates $\gamma$ into the world
frame -- the elevation of link $a$ itself:

\begin{equation}
\raisebox{-0.5\height}{\tikz \node[draw=gray!45!black, line width=1.4pt, rounded corners=3pt,
fill=gray!8, inner sep=14pt]
{\Large $\displaystyle \texttt{shoulder\_angle} = \gamma + \varepsilon$};}
\qquad
\texttt{elbow\_angle} = \delta
\end{equation}

\section{Notation $\leftrightarrow$ Code}

\begin{center}
\begin{tabular}{@{}ll@{}}
\toprule
Symbol & Variable \\
\midrule
$\mathbf s$ & \texttt{shoulder\_servo\_position} \\
$\mathbf w$ & \texttt{wrist\_position} \\
$a,\,b$ & \texttt{shoulder\_arm\_length}, \texttt{elbow\_arm\_length} \\
$\mathbf d$ & \texttt{shoulder\_2\_wrist\_vec} \\
$r$ & \texttt{xy\_projection\_dist} \\
$\varepsilon$ & \texttt{shoulder\_2\_wrist\_vec\_angle} \\
$c$ & \texttt{shoulder\_2\_wrist\_length} \\
$\gamma$ & (\texttt{shoulder\_angle}, before the \texttt{+=} correction) \\
$\delta$ & \texttt{elbow\_angle} \\
$\gamma+\varepsilon$ & \texttt{shoulder\_angle} (returned) \\
\bottomrule
\end{tabular}
\end{center}

\end{document}
